\chapter{Uso de IA}\label{cap:ia}

\section{Declaración de uso}
De acuerdo con las normativas de la asignatura, se declara el uso de herramientas de IA generativa (ChatGPT, Gemini y GitHub Copilot) como asistentes en las fases de ideación, implementación técnica del cliente web y refinamiento de la memoria técnica. Los autores asumen la responsabilidad total del contenido, habiendo supervisado, verificado y modificado cada fragmento de código y referencia bibliográfica integrada en el proyecto.

\section{Herramientas y Propósitos}
El uso de estas herramientas se desglosa según las siguientes necesidades técnicas:

\begin{itemize}
	\item \textbf{ChatGPT / Gemini (Asistencia Conversacional):} Utilizados en la fase inicial para el \textit{brainstorming} de la arquitectura de la dApp. Se debatió la viabilidad de usar un sistema de depósitos \textit{Escrow} frente a intercambios atómicos (\textit{Atomic Swaps}), decidiendo finalmente una implementación híbrida.
	\item \textbf{GitHub Copilot (Modo Plan y Autocompletado):} Empleado como asistente en el desarrollo del frontend. 
	Se utilizó para agilizar la creación de componentes de interfaz y la integración de las APIs externas de Scryfall 
	(Magic: The Gathering) y Pokémon TCG. También se utilizó como asistente en el desarrollo, hizo una primera estructura 
	de la interfaz frontend, facilitó la redacción de scripts personalizados para simplificar el trabajo (como \textit{Start-Local.cmd}, 
	detallado en el Capítulo \ref{cap:manuales}) y ayudó a hacer funcionar MetaMask con nuestra aplicación. 
	\item \textbf{Claude:} Utilizado de forma puntual para la identificación de errores de tipado en TypeScript y para asegurar 
	la cohesión gramatical de los apartados teóricos de la memoria.
\end{itemize}

\section{Metodología y Control de Calidad}
El contenido generado no ha sido incorporado de forma directa. Se ha seguido un proceso de validación 
crítica: \textbf{Prompt $\rightarrow$ Generación $\rightarrow$ Revisión Técnica $\rightarrow$ Modificación Manual $\rightarrow$ Integración $\rightarrow$ Revisión manual}. 

En el desarrollo de código se ha tomado un enfoque de planificación y conversación para mantener las 
ideas claras, mitigar el sesgo de complacencia (sicofancia) del modelo y ceñirnos a los requisitos: \textbf{Prompt (en modo Plan) 
$\rightarrow$ Aclarar dudas del modelo $\rightarrow$ Primera generación $\rightarrow$ Refinar conversando 
$\rightarrow$ Segunda generación (y repetir si es necesario) $\rightarrow$ Revisión Técnica $\rightarrow$ 
Modificación Manual $\rightarrow$ Integración $\rightarrow$ Revisión manual}. 

En el desarrollo del frontend, se realizaron ajustes manuales significativos en la lógica de conexión con \textit{ethers.js} y en la gestión del \textit{Chain ID}. Las herramientas de IA presentaban soluciones genéricas que no contemplaban las particularidades de nuestro entorno de red local en Hardhat, por lo que la intervención humana fue necesaria para asegurar la funcionalidad del \textit{injected provider} siguiendo el estándar EIP-1193. La IA insistía mucho en usar métodos obsoletos como \textit{window.web3} y tuvimos que, manualmente, ajustarlo al estándar moderno \textit{window.ethereum} para gestionar mejor los cambios de red y de cuenta.

\section{Prompts Significativos y Sesiones de Diseño}
A continuación, se documentan las interacciones más críticas que definieron el rumbo del frontend y la arquitectura de datos.

\textit{Nota:} Los prompts estan recortados para no entorpecer la lectura de la memoria. Los originales son mucho 
más largos, estamos en contra de usar la IA de la forma \textit{"hazme esto"}. Le hemos especificado requisitos, 
la hemos dirigido paso a paso, revisado todo lo que realizaba y le hemos pedido, explícitamente, que evite la 
sicofancia o que asienta con todo lo que le decimos. La hemos usado como herramienta, una herramienta muy útil 
que agiliza el trabajo tanto de investigación como de programación y que nos ha permitido no solo no reducir 
el alcance si no de aumentarlo.

\subsection{Sesión de Diseño Inicial (Gemini)}
Enlace a la conversación completa donde planteamos inicialmente la aplicación: \\ \url{https://gemini.google.com/share/add0194ba4c6}

\subsection{Diseño de la Capa de Persistencia y Transaccionalidad (Solidity)}
Instrucción principal para estructurar la lógica de negocio en la cadena de bloques, 
haciendo especial hincapié en la integridad de los datos y la atomicidad, conceptos centrales 
en el diseño de bases de datos robustas:

\begin{quote}
    \textit{"Vamos a modelar la capa de persistencia en Solidity para un mercado de cartas TCG. 
	Mi objetivo principal es ofrecer dos flujos de interacción: compraventa y trueque. Para el flujo de compraventa,
	quiero implementar un sistema de garantía que asegure que el comprador deposite el importe antes de que el vendedor
	decida entregar la carta, y que el vendedor solo pueda reclamar el pago una vez que el comprador haya confirmado la recepción de la carta. 
	En el caso del trueque, quiero que los usuarios puedan proponer intercambios de cartas sin necesidad de intermediarios,
	pudiendo aceptar o rechazar propuestas de intercambio, y que la transferencia de propiedad de ambas cartas se realice de forma atómica.
	A nivel de estructura de 
	datos, ¿cuál sería la forma más eficiente de modelar para minimizar el coste de gas, teniendo en cuenta que cada carta es un NFT con metadatos almacenados en IPFS?	"}
\end{quote}

\subsection{Definición de Requisitos del Frontend (Copilot Plan Mode)}

Este prompt fue importante para establecer el alcance de la interfaz y la integración con la base de datos distribuida e IPFS:

\begin{quote}
	\textit{"Vamos a analizar un proyecto de intercambio y compraventa de cartas TCG usando blockchain. 
	Requisitos: Frontend en React/TS, login mediante MetaMask, detección automática de activos del usuario. 
	Debe incluir creación manual de cartas (Minting) y buscador mediante APIs de Scryfall y Pokémon. 
	Cada carta requiere: Nombre, Imagen y un número de serie PSA generado aleatoriamente. Implementar 
	vista de mercado con historial de intercambios para verificar la trazabilidad en la blockchain. Uso 
	de IPFS mediante el SDK de Pinata para la persistencia off-chain."}
\end{quote}

\subsection{Refinamiento de UI y Consistencia de Datos}
Prompt utilizado para corregir problemas de legibilidad y estados de conexión:

\begin{quote}
	\textit{"Vamos a ajustar la interfaz: en Magic: The Gathering, permitir seleccionar versiones de la carta (estilo mtgprint). En 'Mi Colección', corregir el solapamiento de texto sobre las imágenes de los NFTs. Añadir un indicador visual de estado (Tick verde / X roja) para verificar la sincronización entre el Chain ID de la wallet y el Chain ID de la red Hardhat (31337)."}
\end{quote}